\section{Đánh giá hạn chế, giải pháp cải tiến và định hướng hoàn thiện}

\subsection{Hạn chế và giải pháp tối ưu ở giai đoạn thu thập dữ liệu}

Phân tích nhật ký vận hành thực tế trên toàn bộ kho $25{,}000$ bài báo cho thấy điểm nghẽn dễ nhận thấy nhất ở khâu phát hiện bài báo là sự lãng phí yêu cầu mạng tại các dải dữ liệu vắt qua ranh giới tháng. Để thu được đúng $25{,}000$ bài báo hợp lệ, hệ thống đã phải dò tổng cộng $179{,}012$ mã ứng viên trên toàn nhóm. Trong khi ba dải nằm trọn trong một tháng chỉ tốn đúng $5{,}000$ lượt dò mỗi dải ($168.7$ đến $198.9\text{ s}$), hai dải xuyên tháng của Phạm Minh Trọng (tháng $04$--$05/2025$) và Nguyễn Phúc Khang (tháng $06$--$07/2025$) lần lượt tiêu tốn tới $83{,}142$ lượt dò ($3{,}276.5\text{ s}$) và $80{,}870$ lượt dò ($2{,}721.8\text{ s}$). Nguyên nhân trực tiếp là dãy số thứ tự hàng tháng của arXiv luôn dừng thấp hơn rất nhiều so với trần năm chữ số $99{,}999$, khiến chiến lược quét tuần tự phải đi qua đúng $154{,}012$ mã trống ở phần đuôi tháng trước khi chuyển sang tháng kế tiếp. Giải pháp kỹ thuật triệt để là thay cơ chế quét tuyến tính bằng bước ước lượng ranh giới tháng hai pha gồm nhảy bước lũy thừa kết hợp tìm kiếm nhị phân để xác định chính xác chỉ số bài báo cuối cùng hợp lệ trong tháng đầu tiên trước khi gửi truy vấn thu hoạch siêu dữ liệu theo lô qua giao thức OAI-PMH~\cite{arxiv_api,arxiv_oai_2025}. Cách tiếp cận này hạ số yêu cầu thăm dò ranh giới từ hàng chục nghìn xuống mức $O(\log K)$ (dưới hai mươi lượt kiểm tra tồn tại cho mỗi giao điểm tháng), rút ngắn thời gian dò tìm trên các dải xuyên tháng về ngang bằng các dải nội tháng mà vẫn giữ tính tất định và tuân thủ nghiêm giới hạn tần suất truy cập máy chủ.

Điểm nghẽn thời gian thực thi lớn thứ hai nằm ở công đoạn thu hoạch $466{,}763$ liên kết tài liệu tham khảo arXiv cho $20{,}878$ bài báo trích dẫn từ đồ thị tri thức Semantic Scholar~\cite{kinney2023,ammar2018}, vốn chiếm tỷ trọng chủ đạo trong tổng thời gian chạy $26.02$ đến $35.48\text{ giờ}$ trên Google Colab của từng thành viên. Việc gọi từng bài báo đơn lẻ bằng yêu cầu GET chịu giới hạn tần suất nghiêm ngặt tạo ra độ trễ chờ cộng dồn lớn, đồng thời gây truy vấn lặp lại đối với những công trình nền tảng được trích dẫn hàng trăm lần trong cùng kho $25{,}000$ bài. Chuyển sang kiến trúc thu hoạch theo lô thích ứng qua phương thức POST kết hợp bộ điều tiết thùng thẻ bài động, chỉ quay về truy vấn đơn lẻ cho các bài báo có danh sách tham khảo quá dài bị cắt trang, cùng bộ nhớ đệm cục bộ ánh xạ siêu dữ liệu các mã tham chiếu đã giải mã sẽ cắt giảm đáng kể số vòng gọi mạng. Song song với đó, cơ chế quản lý trạng thái tải mã nguồn cần tách bạch hoàn toàn giữa lỗi định dạng bất khả kháng (chẳng hạn $5.38\%$ đến $7.28\%$ phiên bản chỉ có tệp PDF hoặc định dạng không hỗ trợ, nhất quán với khảo sát quy mô lớn~\cite{ginev2024}) với lỗi tạm thời do giới hạn lưu lượng (HTTP 429) hay sự cố máy chủ (HTTP 5xx) sau khi cạn hạn mức thử lại tức thời. Thay vì chốt sớm phiên bản hết lượt thử lại vào trạng thái không thể xử lý vĩnh viễn, cơ sở dữ liệu điểm phục hồi quan hệ cần đánh dấu rõ cờ có thể thử lại để các lượt quét bù độ trễ dài tự động thu hồi những phiên bản gặp nghẽn mạng nhất thời mà không tải lại các bản không có mã nguồn LaTeX.

Về mặt đo lường hiệu năng và quản trị thực nghiệm, hai khoảng trống quan sát cần được chuẩn hóa để nâng cao tính minh bạch tái lập giữa các phiên chạy. Thứ nhất, số liệu đo trên $167$ phiên chạy Colab tại mốc ngày $06/11/2026$ ghi nhận thời gian xử lý trung bình mỗi bài báo ($10.64$ đến $22.61\text{ s}$) lớn gần gấp đôi trung vị ($4.71$ đến $9.56\text{ s}$). Hiện tượng phình số đo này phát sinh khi phép tính lấy mốc hoàn thành cuối cùng trừ mốc bắt đầu đầu tiên vô tình gộp cả khoảng thời gian nằm chờ trong hàng đợi khi bộ điều phối đa giai đoạn chạy xen kẽ nhiều bài báo. Để tách bạch chi phí tính toán và mạng thực tế khỏi độ trễ xếp hàng, hệ thống đo cần ghi nhận độc lập thời lượng hoạt động của từng công đoạn trên mỗi bài báo rồi cộng dồn làm thời gian thực thi hữu ích, đồng thời báo cáo riêng thời gian lưu trú trong hàng đợi. Thứ hai, các bản ghi chuẩn mốc lịch sử ngày $06/11/2026$ được thu trước khi hệ thống bổ sung xác nhận tự động môi trường không dùng GPU và trước đợt bù đắp tài liệu tham khảo cuối cùng cho ba dải dữ liệu, trong khi nhật ký đo thời gian toàn cục của Giai đoạn 2 vẫn lưu phân tán trên bộ nhớ đám mây ngoài kho lưu trữ nội bộ. Hợp nhất cơ chế ghi nhật ký có cấu trúc cho cả Giai đoạn 1 và Giai đoạn 2 sau bước bù đắp cuối cùng---kèm chứng thực phần cứng CPU, phân rã thời lượng từng công đoạn và chuỗi mẫu bộ nhớ tiến trình theo thời gian---sẽ khép kín chuỗi kiểm toán hiệu năng đầu-cuối trên toàn kho $25{,}000$ bài báo.

\subsection{Hạn chế và giải pháp nâng cao độ chính xác ở giai đoạn phân tích cú pháp}

Quá trình phân tích cú pháp trên $22{,}961$ bài báo ($37{,}967$ phiên bản, đạt $92.84\%$ tổng số $40{,}896$ phiên bản khai báo) đã chuyển đổi $16{,}223{,}222$ lượt xuất hiện phần tử thành $11{,}098{,}189$ phần tử duy nhất (giảm $31.59\%$ nhờ khử trùng lặp liên phiên bản) cùng $16{,}185{,}255$ cạnh quan hệ phân cấp. Tuy nhiên, kiểm tra định tính trên đầu ra cho thấy quy trình chuẩn hóa biểu thức toán học hiện tập trung chủ yếu vào các đoạn văn xuôi thân bài trước khi tách câu, còn công thức nội tuyến nằm trong chú thích hình bảng, ô bảng số liệu và tiêu đề mục lục chưa được chuẩn hóa đồng đều theo cùng một chuẩn. Sự bất đối xứng này khiến những chỉnh sửa nhỏ về cú pháp toán học trong tiêu đề mục hoặc chú thích hình giữa hai phiên bản vẫn làm thay đổi chuỗi văn bản và bỏ lỡ cơ hội khử trùng lặp, đồng thời giảm độ nhất quán khi tìm kiếm biểu thức toán trên $285{,}123$ phần tử hình bảng và $468{,}790$ tiêu đề phân cấp (gồm $184{,}592$ mục lớn, $222{,}075$ tiểu mục và $62{,}123$ tiểu mục cấp ba). Mở rộng đường ống chuẩn hóa toán học áp dụng đồng nhất cho mọi nút văn bản thuộc tiêu đề, chú thích và ô bảng là cải tiến có chi phí tính toán thấp nhưng mang lại hiệu quả trực tiếp cho chất lượng chỉ mục. Song song đó, tầng nhận diện cấu trúc cần mở rộng phạm vi bao phủ sang các môi trường bảng có độ rộng tự co giãn, hình bảng xoay ngang và danh sách mô tả định nghĩa. Việc xử lý tường minh cả các khối công thức đặt ngay sau lệnh ngắt dòng mà không có dòng trống ngăn cách cũng sẽ giúp phân tách sạch $991{,}526$ công thức khối khỏi các câu văn xuôi liền kề.

Ở tầng tách câu, hệ thống sản xuất cho $8{,}634{,}050$ câu duy nhất lựa chọn bộ tách theo quy tắc tất định (đạt $7/12$ trường hợp chẩn đoán khó và $40/40$ bài kiểm tra tương thích) nhằm đảm bảo khởi chạy tức thời không phụ thuộc trọng số ngoài và tái lập hoàn toàn theo từng điểm phục hồi trên Colab. Tuy nhiên, cấu trúc biểu thức chính quy kèm danh mục hai mươi từ viết tắt vẫn phát sinh lỗi cắt câu sớm khi gặp chữ cái viết tắt tên tác giả ngắn hoặc có dấu gạch nối, các chuỗi viết tắt học thuật nhiều dấu chấm và cụm trích dẫn dày đặc xen giữa mệnh đề. Ngay cả phương pháp không giám sát NLTK Punkt~\cite{kiss2006} cũng chỉ nhích lên $8/12$ trường hợp chẩn đoán nhưng lại làm giảm tỷ lệ tương thích xuống $38/40$. Thực nghiệm đo chuẩn có kiểm soát trên $100$ bài báo đầu dải ($93$ bài có mã nguồn LaTeX, $7{,}525$ khối văn xuôi, $4{,}803{,}251$ ký tự) mở ra hướng nâng cấp khả thi rõ rệt: mô hình tách câu thống kê nhẹ của spaCy~\cite{spacy} đạt độ chính xác chẩn đoán cao nhất ($10/12$) và khi chạy theo lô kích thước $32$ chỉ tốn $0.104\text{ s}$ mỗi phiên bản với $163.6\text{ MiB}$ bộ nhớ đỉnh (tốc độ $495{,}000$ ký tự mỗi giây). Mức thông lượng này ngoại suy chỉ mất khoảng $1.16\text{ giờ}$ cho toàn bộ $40{,}000$ phiên bản LaTeX, tiết kiệm hơn hai bậc độ lớn so với mô hình nơ-ron đa ngôn ngữ SaT~\cite{frohmann2024} trên CPU đơn lẻ ($153.39\text{ giờ}$ và $1{,}310.1\text{ MiB}$). Do đó, triển khai spaCy chạy theo lô làm bộ đề xuất ranh giới chính kết hợp tầng hậu kiểm tất định cấm cắt ngang bên trong công thức nội tuyến, ngoặc trích dẫn và mẫu viết tắt tên tác giả sẽ đồng thời nâng độ chính xác ngôn ngữ và bảo toàn trọn vẹn tính toàn vẹn cấu trúc toán học trên phần cứng CPU phổ thông.

Một nguồn sai số ngược dòng có tác động trực tiếp đến Giai đoạn 3 nằm ở bước trích xuất mục tham khảo nội tuyến và gộp thư mục BibTeX liên phiên bản. Toàn bộ kho dữ liệu hiện gom được $975{,}674$ mục BibTeX từ $19{,}204$ bài báo ($76.82\%$ tổng số bài) với độ phủ tiêu đề $99.23\%$, năm $98.36\%$ và tác giả $97.03\%$, nhưng vẫn còn $19.17\%$ mục mang loại tổng quát dự phòng thay vì phân loại cụ thể như bài báo tạp chí ($57.93\%$) hay kỷ yếu hội nghị ($18.25\%$). Như phân tích lỗi tại Bảng~\ref{tab:m3_error_forensics} đã chỉ rõ ở trường hợp bất đồng thứ năm, các mục tham khảo nội tuyến trong thiên văn học dùng ký hiệu gạch ngang thay thế tên tác giả lặp lại và nhúng mã arXiv vào trong chuỗi DOI tạp chí khiến bộ phân tích mục nội tuyến hiện tại trích xuất sai trường tác giả và làm nhiễu chuỗi tiêu đề. Để khắc phục triệt để, khối xử lý thư mục cần bổ sung bộ giải tham chiếu tuần tự cho ký hiệu gạch ngang lặp tác giả bằng cách kế thừa danh sách tác giả đã chuẩn hóa từ mục tham khảo liền trước. Tiếp theo, quy tắc gộp liên phiên bản cần áp dụng thứ tự ưu tiên phân loại để loại bài báo tạp chí hoặc kỷ yếu hội nghị tự động thay thế loại dự phòng tổng quát mỗi khi một phiên bản khác bổ sung thông tin tạp chí hay hội nghị (hiện đạt $71.11\%$ độ phủ trường tạp chí). Cuối cùng, bổ sung cơ chế khử trùng lặp mờ đối với biến thể cách viết tên tác giả giữa dạng viết tắt chữ cái đầu và tên đầy đủ sẽ nâng cao tính nhất quán của siêu dữ liệu tác giả cho toàn bộ các tác vụ hạ nguồn.

\subsection{Hạn chế và giải pháp hoàn thiện mô hình đối sánh tài liệu tham khảo}

Mặc dù mô hình Hồi quy Logistic đạt MRR@5 tuyệt đối $1.0000$ trên cả hai tập kiểm định độc lập và chỉ mắc đúng $1$ sai sót xếp hạng thực chất trên toàn bộ $35{,}373$ truy vấn, các kết quả kiểm toán ngược và phân tích lỗi đã chỉ rõ những nút thắt cần xử lý khi mở rộng quy mô. Trước hết, đối với bộ quy tắc gán nhãn tự động tạo dữ liệu giám sát yếu (đạt độ chính xác $98.72\%$ và độ phủ $98.83\%$ trên $25$ bài báo thủ công), toàn bộ $24/24$ lỗi dương tính giả đều tập trung ở tầng mờ thứ ba và $21/22$ lỗi bỏ sót xuất phát từ cơ chế phủ quyết cứng khi vắng mã arXiv. Cải tiến trọng tâm là siết chặt bộ trích xuất định danh bằng yêu cầu ngữ cảnh arXiv tường minh kết hợp chuẩn hóa các mã chuyên ngành kiểu cũ trước năm $2007$ về dạng thống nhất trước khi so sánh, đồng thời điều chỉnh tầng chuẩn hóa tiêu đề của bộ quy tắc để bảo toàn các ký hiệu toán học ngắn hai ký tự như $C_6$ hay $W_4$. Bên cạnh đó, chuyển cơ chế phủ quyết cứng khi mã arXiv vắng mặt thành một mức phạt mềm có thể vượt qua khi tiêu đề, tác giả và năm công bố đạt đồng thuận rất cao sẽ giảm mạnh tỷ lệ bỏ sót mà không làm tăng lỗi ghép nhầm. Về mặt thiết kế thực nghiệm, nghiên cứu loại trừ đặc trưng trong tương lai cần tính lại các đặc trưng cạnh tranh nhóm hoàn toàn trên điểm tổng hợp của các nhóm còn lại khi loại nhóm Tiêu đề nhằm tháo gỡ sự ghép dính phương pháp luận hiện tại. Song song đó, bổ sung tầng chuẩn hóa biến thể từ ghép có và không có dấu gạch nối (đồng nhất các cặp như ``nonflat'' và ``non-flat'') sẽ xử lý trực tiếp nguyên nhân dẫn tới sai số xếp hạng thực chất duy nhất được ghi nhận tại Bảng~\ref{tab:m3_error_forensics}.

Ở cấp độ mô hình hóa bài toán, kiến trúc xếp hạng độc lập theo từng truy vấn hiện tại vẫn còn hai khoảng trống cấu trúc đáng kể khi triển khai trong môi trường sản xuất thực tế. Khoảng trống thứ nhất là khả năng phát hiện tường minh các mục trích dẫn không có bài báo đích trên arXiv: trên tập $101{,}441$ mục BibTeX của nhóm, có tới $65{,}816$ mục không thuộc kho arXiv đã thu hoạch, nhưng mô hình xếp hạng hiện chỉ sắp thứ tự năm ứng viên tốt nhất mà chưa đưa ra quyết định từ chối có kiểm soát độ tin cậy. Khắc phục giới hạn này đòi hỏi huấn luyện thêm nhánh quyết định có hoặc không khớp trên cả các truy vấn toàn âm kết hợp hiệu chỉnh xác suất (như hiệu chỉnh Platt hay hồi quy đơn điệu) để đặt ngưỡng từ chối tối ưu theo tỷ lệ chính xác--độ phủ mong muốn. Khoảng trống thứ hai là tính độc lập giữa các truy vấn trong cùng một bài báo trích dẫn, vốn cho phép hai mục BibTeX khác nhau cùng giành một tài liệu tham khảo arXiv duy nhất làm dự đoán hạng nhất. Áp dụng thuật toán ghép cặp hai phía trọng số cực đại toàn cục (thuật toán Hungary) lên ma trận điểm $m_p \times n_p$ đã hiệu chỉnh của từng bài báo sẽ cưỡng chế ràng buộc tương ứng một-một tự nhiên và tận dụng thông tin loại trừ lẫn nhau. Đáng chú ý hơn, do Giai đoạn 2 đã lưu giữ đầy đủ cấu trúc cây văn bản gắn liền với từng khóa trích dẫn, việc bổ sung kênh đặc trưng ngữ cảnh trích dẫn cục bộ---đo độ tương đồng từ vựng giữa các câu văn bao quanh vị trí trích dẫn trong bài báo với tiêu đề và tóm tắt của tài liệu ứng viên---sẽ cung cấp một nguồn tín hiệu trực giao hoàn toàn mới khi siêu dữ liệu thư mục bị thiếu hoặc cắt xén mà vẫn tuân thủ giới hạn không phụ thuộc mô hình ngôn ngữ lớn bên ngoài.

Kết quả thí nghiệm che mã định danh tại Bảng~\ref{tab:m3_results} và công tác đóng gói phát hành cũng gợi mở lộ trình nâng cấp rõ ràng cho cả lõi mô hình lẫn hạ tầng dữ liệu. Khi áp dụng cho các kho học thuật cũ hơn hoặc các chuyên ngành ít gắn mã arXiv trực tiếp trong tệp thư mục, mô hình xếp hạng theo danh sách LightGBM LambdaRank ($31$ lá) kết hợp tăng cường dữ liệu che $50\%$ mã định danh thể hiện độ bền rõ rệt trước sự vắng mặt của mã định danh. Cụ thể, cấu hình này đã cắt giảm khoảng $70\%$ sai số kiểm định chéo khi che mã (từ $0.9970$ lên $0.9991$) trong khi vẫn giữ trọn mức $1.0000$ trên cả hai tập kiểm định và kiểm tra độc lập. Việc kết hợp có trọng số giữa xác suất đã hiệu chỉnh của Hồi quy Logistic với điểm xếp hạng của LightGBM LambdaRank tăng cường che mã có thể vận hành gọn nhẹ trong cùng giới hạn tài nguyên CPU hiện tại ($12.88\text{ phút}$ chạy toàn nhóm, bộ nhớ đỉnh $2{,}393.5\text{ MB}$ trên $2$ luồng CPU). Về mặt quản trị dữ liệu mở, dù năm bộ dữ liệu cột Parquet phân mảnh theo từng thành viên (tổng dung lượng $1{,}199.45\text{ MiB}$, tương đương $1.171\text{ GiB}$ trên sáu bảng chuẩn) đều đã vượt qua quy trình kiểm định lược đồ, khóa chính, khóa ngoại và quy tắc xếp hạng với $0$ lỗi, kho dữ liệu hiện vẫn phân tán theo năm đầu mục riêng lẻ. Hợp nhất năm kho thành viên dưới một thẻ dữ liệu chung duy nhất của toàn nhóm---đi kèm bảng kiểm kê mã kiểm tra toàn vẹn, chỉ số trùng lặp liên thành viên và tập truy vấn đối chuẩn tái lập---sẽ chuyển hóa năm sản phẩm học phần song song thành một kho chuẩn thống nhất, sẵn sàng phục vụ nghiên cứu phân tích mạng trích dẫn khoa học.
